\section{Architecture}
\label{sec:architecture}

This section describes the architecture of the \tradesq{} Mobile Application using the C4 model.
The C4 model describes the system at increasing levels of technical detail through System Context, Container, and Component views.
The C1 view describes the application within its external environment.
The C2 view describes the major runtime and local storage containers of the mobile application.
The C3 view describes the internal components of the Expo React Native Application container.

A C4 code-level section is omitted because code-level diagrams become outdated more quickly than the higher-level architecture.
The C3 component diagram provides sufficient detail to explain the application's main implementation structure.

\subsection{System Context (C1)}
\label{sec:c1}

The \tradesq{} Mobile Application is the user-facing part of \tradesq{}, a trade trading application intended to help people exchange jobs and establish social connections within their local neighbourhood.
The primary target group includes retired older adults.
Users sign in before accessing application functionality.
Authenticated users can create trade offers and requests, discover listings, manage their profile, save favourites, express interest in listings, upload images, and view notifications.

\begin{figure}[H]
    \centering
    \includegraphics[width=\textwidth]{hd/diagrams/c1}
    \caption{C1 System Context diagram of the \tradesq{} Mobile Application.}
    \label{fig:c1}
\end{figure}

As shown in \cref{fig:c1}, the \tradesq{} User interacts directly with the \tradesq{} Mobile Application.
The application provides the user interface through which authenticated users access trade listings, profiles, favourites, conversations, and notifications.

The \tradesq{} Back-end API is an external dependency of the mobile application.
The application communicates with the API through HTTPS to retrieve and update shared application data, including authentication, profiles, listings, favourites, notifications, and other user-facing information.
Notifications are retrieved through authenticated HTTPS requests and displayed in the application notification interface.
The application does not use WebSocket communication because real-time messaging is outside the current application scope.

The application does not directly access server-side persistence or storage systems.
Those systems are implementation details of the Back-end and are outside the scope of this document.

The application also uses services offered by the mobile operating system.
These platform services provide secure storage for session credentials, network-connectivity information, media selection for profile and listing images, notification capabilities, and accessibility functionality such as screen readers and system text scaling.
